\section*{Chapter 2: Background}
\setcounter{section}{2} \setcounter{subsection}{0}
\phantomsection \label{chap:background} \addcontentsline{toc}{section}{\textit{Chapter 2: Background}}



\subsection{Agent Systems and the Limits of User-space Control Planes}

Modern operating systems are built around the process as the primary unit of protection: address-space isolation, controlled file access, and well-defined IPC \cite{tanenbaum2022}.
These abstractions remain fundamental, but agent-based workloads introduce a layer of intent between the user and execution that they do not directly address.
An agent plans over semantically described tools, constructs structured requests, and acts on behalf of a user; correctness now depends not only on process isolation but on whether the semantics of a requested operation match what is actually executed.

Several lines of work have responded to this gap.
Systems such as generative agents demonstrate long-lived entities capable of maintaining memory and coordinating workflows \cite{park2023generative}.
Frameworks including LangChain \cite{langchain2022}, LangGraph \cite{langgraph2024}, AutoGPT \cite{autogpt2023}, CrewAI \cite{crewai2024}, and Semantic Kernel \cite{semantickernel2023} structure tool composition in multi-step pipelines.
More ambitious proposals embed LLM-centric services directly into a system-level kernel managing scheduling, memory, and tool access \cite{packer2023memgpt, wu2024oscopilot, mei2024aios, liu2026agentos}.
At the application layer, \texttt{OpenClaw} \cite{openclaw2026} acts as a persistent user-space control layer for tool execution; Harness \cite{harness2025} embeds governed AI agents in CI/CD pipelines; Apple Intelligence \cite{apple2024pcc} integrates agents with on-device OS services and a private-cloud compute layer.

Section~\ref{sec:limits-user-space} details why this enforcement placement matters.



\subsection{What MCP Provides}

The Model Context Protocol (MCP), released by Anthropic in November 2024, provides a useful abstraction for structuring tool-mediated agent interaction \cite{mcp2025}.
It defines a client-host-server architecture in which tools, resources, and prompts are exposed through a standardised JSON-RPC interface.
The host coordinates sessions and mediates access, while servers implement capabilities; tools are the primary mechanism by which agents take executable actions on external systems.
MCP rapidly achieved broad ecosystem adoption, with implementations in Python, TypeScript, C\#, and Java, and integration across major LLM providers and agent frameworks \cite{mcp2025}.

This separation simplifies interoperability and modularity.
However, it also concentrates responsibility in the host: identity management, approval decisions, and consistency between tool descriptions and execution all depend on correct host behaviour.
MCP defines how components communicate, but does not specify how a local system should enforce these properties.
\texttt{linux-mcp} implements the enforcement layer that MCP leaves unspecified: it provides the local mechanism that turns the protocol's abstract trust model into a set of kernel-verified invariants.

\subsection{Limits of User-space Mediation}
\label{sec:limits-user-space}

Chapter~\ref{chap:introduction} introduced three control-plane invariants; this section explains their technical origin at the MCP protocol layer.

\textbf{Semantic drift.}
Tool invocation depends on semantically rich manifests combining natural-language descriptions, schemas, and examples \cite{schick2023toolformer, patil2023gorilla, mcp2025}.
Because both interpretation and enforcement reside in a mutable user-space process, a change in implementation or a weakened control path may result in executing a capability that no longer matches the approved semantics \cite{liu2023prompt}.

\textbf{Identity anchoring.}
Agent systems bind requests to sessions or agents rather than PIDs \cite{yao2023react, autogpt2023, park2023generative}, but MCP provides no mechanism to anchor these identities in OS-visible credentials.
Prompt injection and session confusion attacks can therefore cause agents to exceed their intended authority even when the gateway appears to enforce correct behaviour \cite{greshake2023prompt, openclaw_security1, openclaw_security2}.

\textbf{Durability.}
Approval decisions and pending requests may be lost on crash or restart when control-plane state is held only in memory \cite{packer2023memgpt}.
The fragility is particularly acute in long-running agent systems where execution spans multiple steps and decisions \cite{wang2023voyager, yang2024swe}.

These limitations are not unique to any single framework; they reflect a broader pattern in which control-plane invariants \cite{lim2024compartmentalization} are difficult to enforce when confined to user space \cite{schneider2022sok, lefeuvre2024sok, ji2026seagent, shi2025progent}.




\subsection{Relation to Native Linux Security Mechanisms}

Any such design must coexist with existing Linux security mechanisms.
Facilities such as seccomp, AppArmor, SELinux, and Landlock provide strong guarantees over system calls, resources, and process privileges \cite{seccomp2026, lsm2026, apparmor2026, selinux2026, landlock2026}.
These mechanisms operate at the level of kernel objects and access rights.

However, none of them capture higher-level concepts such as tool identity, semantic equivalence between a manifest and an execution endpoint, or the scope of an approval decision.
The root cause is a layer mismatch: existing mechanisms enforce policy at the level of kernel objects (files, processes, syscalls), whereas tool identity is an application-layer concept defined by a semantic hash of a JSON manifest.
Extending these mechanisms to enforce tool-identity would require either encoding tool hashes into LSM policy (conflicting with the goal of protocol-layer independence) or hooking into the application-protocol layer, which breaks the principle that LSM hooks attach to kernel operations, not to application messages.

This gap motivates introducing a dedicated control layer that validates semantic and identity invariants at the point of tool execution, outside the mutable execution path of the user-space gateway.



\subsection{Related Work}
\label{sec:related-work}

Several existing mechanisms address related but distinct problems; understanding why they fall short for MCP-style enforcement motivates the present design.

\vspace{4pt}
\noindent\textit{Kernel-native enforcement mechanisms.}

\textbf{Seccomp-BPF.}
Seccomp-BPF can filter syscalls and inspect syscall arguments at the process boundary~\cite{seccomp2026}, but tool identity is not expressed at the syscall level; it is embedded in the JSON-RPC payload exchanged over Unix domain sockets.
Seccomp sees only the \texttt{socket()} and \texttt{sendto()} syscalls; it cannot inspect the message content to verify that a tool invocation matches an approved semantic hash or that an approval ticket is unconsumed.
This is a layering issue: tool identity is a concept at the protocol layer, not the syscall layer.
Applied to \texttt{mcpd}, seccomp can reduce the syscall surface but cannot enforce the invariants on which admission control depends.

\textbf{eBPF and LSM hooks.}
eBPF programs can enforce byte-level socket policies with very low overhead \cite{lsm2026}, and LSM hooks attach to VFS and socket operations at well-defined kernel hook points \cite{brimhall2023mac}.
Neither, however, can natively maintain a dynamic, per-tool hash registry and perform variable-length hash comparisons without significant complexity: the BPF verifier's loop and memory restrictions make runtime registry lookups impractical, and LSM hooks fire at the system-call boundary rather than at the application-protocol boundary where tool identity is meaningful (see Section~\ref{sec:design-alternatives} for a more detailed analysis).

\textbf{Landlock.}
Landlock provides process-scoped filesystem sandboxing through unprivileged LSM restrictions \cite{landlock2026}.
It is complementary to this work (a tool service could be Landlock-sandboxed to limit its filesystem reach) but does not address tool identity, approval binding, or approval-state durability across daemon failures.

\textbf{Capsicum.}
Capability-based designs scope authority to individual file descriptors or object references \cite{lefeuvre2024sok}.
Adapting this to MCP-style tool invocation would require tool services to be written against a capability API, violating the design goal that tool authors need not be aware of any kernel interface.
More fundamentally, per-descriptor capabilities cannot natively encode whether a specific tool invocation matches an approved semantic description.

These mechanisms operate below the protocol layer and cannot natively represent tool identity, which is defined by a semantic hash over a JSON manifest.

\vspace{4pt}
\noindent\textit{Agent runtimes and frameworks.}

\textbf{Agent frameworks.}
LangChain \cite{langchain2022}, LangGraph \cite{langgraph2024}, AutoGPT \cite{autogpt2023}, CrewAI \cite{crewai2024}, Haystack \cite{haystack2023}, and Semantic Kernel \cite{semantickernel2023} all locate the full control plane in user space.
They provide no mechanism for durable, kernel-verified approval state; correctness depends entirely on the in-process runtime.
\texttt{OpenClaw} \cite{openclaw2026} operationalises this pattern at scale, and independent security analyses demonstrate the resulting vulnerabilities \cite{openclaw_security1, openclaw_security2}.
Industrial CI/CD agents such as Harness \cite{harness2025} add governance through platform RBAC, but their trust anchor remains the cloud control plane rather than any local kernel boundary.

\textbf{MAC frameworks for LLM agents.}
Recent work has begun applying mandatory access control (MAC) concepts to LLM agent privilege escalation \cite{ji2026seagent} and to programmable privilege control for tool-calling agents \cite{shi2025progent}.
Both represent application-layer MAC: enforcement lives inside a modified agent runtime, so its correctness depends on the runtime itself remaining uncompromised.
\texttt{linux-mcp} moves enforcement out of the gateway entirely into a kernel module, so its invariants hold even if the gateway is buggy or partially compromised.
The two approaches are complementary rather than competing: a runtime-MAC agent framework could run on top of \texttt{linux-mcp}, gaining kernel-level durability without replacing its own semantic enforcement.
The principle of separating policy from enforcement mechanism at the kernel level is well established in Linux MAC frameworks \cite{brimhall2023mac}; \texttt{linux-mcp} applies the same principle to MCP-style tool admission rather than to traditional file and process access control.

All of these keep the trust anchor in user space; their correctness degrades when the runtime itself is compromised.

\vspace{4pt}
\noindent\textit{AIOS-style OS proposals.}

Several recent systems propose embedding LLM-specific abstractions into a new OS substrate \cite{wu2024oscopilot, packer2023memgpt, mei2024aios, liu2026agentos}.
These differ from this work in scope: they aim to replace or fundamentally extend the operating system for agent workloads, managing scheduling, memory, and tool access at the OS level.
\texttt{linux-mcp} takes the opposite approach, leaving the existing OS unchanged and adding only a narrow kernel module for the invariants that user space cannot reliably hold.

These proposals broaden scope by redesigning OS substrates; \texttt{linux-mcp} instead adds a narrow kernel role to an unchanged Linux.

\vspace{4pt}
\noindent\textit{Summary.}

Existing Linux mechanisms operate below the protocol layer and cannot express tool identity.
Existing agent frameworks keep control-plane trust in user space.
AIOS-style proposals broaden scope by redesigning the OS substrate.
\texttt{linux-mcp} instead introduces a narrow kernel role for a bounded set of MCP-style control-plane invariants, enforcing all three of the weaknesses identified in Section~\ref{sec:limits-user-space} outside the gateway's own privilege domain; Chapter~\ref{chap:design} describes how it does so with a single, narrow kernel addition.


\begin{table}[t]
\centering
\small
\begin{tabular}{l p{4.8cm} p{5.2cm}}
\toprule
\textbf{Attack class} & \textbf{Representative attack} & \textbf{Defense in \texttt{linux-mcp}} \\
\midrule

Spoofing
& Forged session identity or mismatched tool hash
& Kernel-verified identity binding and semantic-hash consistency \\

Replay
& Reuse of stale approval tickets or prior requests
& Kernel-enforced approval lifecycle with binding and lifetime validation \\

Substitution
& Executing a tool diverging from declared semantics
& Verification of manifest-derived semantic identity \\

Escalation
& Bypassing gateway checks via logic flaws
& Mandatory kernel arbitration for all invocations \\

\bottomrule
\end{tabular}
\caption{Control-plane threats and defences (in scope).}
\label{tab:threat_in_scope}
\end{table}


\subsection{Threat Model}

The above limitations suggest that the correctness of MCP-style systems depends on a small set of control-plane invariants, including identity binding, approval validity, and consistency between tool semantics and execution.
We therefore ground our design in a threat model that focuses on attacks against these invariants in a local-host setting.

The local-adversary model is appropriate for MCP-style deployments for three reasons.
First, tool services in a typical deployment run on the same host as the gateway; the attack surface that most immediately needs hardening is therefore local, not remote.
Second, a remote adversary who has already achieved code execution on the host is, by definition, a local adversary: the local threat model subsumes this class without requiring a separate treatment.
Third, the two dominant remote-oriented threats in AI agent systems (prompt injection and supply-chain compromise of manifests) operate at different layers and fall outside the scope of kernel-level admission control, as discussed later in this section.
The threat model therefore focuses on what kernel enforcement can actually address: a co-located process that can observe, forge, or replay control-plane messages.

Concretely, the adversary model has the following capabilities: (1) arbitrary code execution on the host as an unprivileged process; (2) access to Unix domain sockets exposed by the gateway; (3) ability to observe, craft, and replay JSON-RPC messages; and (4) full knowledge of the MCP protocol, tool manifests, and \texttt{linux-mcp} implementation (consistent with Kerckhoffs's principle).
The attacker interacts exclusively through the \texttt{mcpd} JSON-RPC surface and holds no elevated capabilities: specifically, no \texttt{CAP\_NET\_ADMIN} (required to invoke the kernel module via Netlink), no \texttt{ptrace} access to \texttt{mcpd}, and no write access to manifest files.
These three boundaries are critical: \texttt{CAP\_NET\_ADMIN} is the interface to kernel enforcement; \texttt{ptrace} would allow process manipulation; manifest write would break the trusted-input assumption.

Two dominant attack categories in AI agent systems fall outside this scope.
\emph{Prompt injection} operates at the model's reasoning layer, before any \texttt{tool:exec} request is submitted; kernel admission control has no visibility into the planning process \cite{greshake2023prompt, liu2023prompt}.
\emph{Supply-chain compromise of manifests} requires file-system-level defences such as IMA or dm-verity; \texttt{linux-mcp} treats on-disk manifests as trusted input and assumes the adversary lacks write access to \texttt{tool-app/manifests/}.
Also out of scope: process-level attacks on the gateway (ptrace, memory corruption, LD\_PRELOAD), kernel-level attacks (rootkits, hardware side-channels), and remote adversaries.

Within this scope, we focus on four classes of control-plane failures:
(1) \emph{spoofing} of identity or session context,
(2) \emph{replay} of previously valid requests or approval decisions,
(3) \emph{substitution} of tool identity or semantics, and
(4) \emph{escalation} through bypassing admission checks in the gateway.
These attack classes and their corresponding defences are summarised in Table~\ref{tab:threat_in_scope}.

These four attack classes directly inform the design of \texttt{linux-mcp} and are revisited in Chapter~\ref{chap:evaluation} through targeted adversarial experiments.

